\section{Executive Summary}

The project is centered on \textbf{retrieval-augmented reward modeling} for
\textit{long-horizon robot manipulation} in simulation. The core idea is simple:
instead of letting a policy rely only on \textit{sparse reward} or hand-crafted shaping,
the system retrieves relevant evidence from a multimodal \textit{external memory}
containing video demonstrations, stage descriptions, instructions, and manual text.
The retrieved context is then fed into a reward model that produces
\textit{dense progress signals} or stage-aware scores, which can in turn support
offline RL, policy selection, or trajectory evaluation.

The proposal is careful not to sell this as a vague ``Embodied RAG'' slogan,
because both retrieval for embodied agents and reward modeling for robotics already
exist as active research directions. Instead, the intended contribution is narrower and easier to evaluate:
\begin{itemize}
  \item use multi-source retrieval to improve \textbf{reward quality} in long-horizon manipulation;
  \item focus on \textbf{OOD states}, \textbf{stage-aware feedback}, and \textbf{sub-trajectory retrieval};
  \item evaluate on \textbf{standard simulation benchmarks} rather than claiming a full real-robot system.
\end{itemize}

\section{Background and Motivation}

In robot manipulation, especially for multi-stage tasks such as opening a drawer and then placing an object,
sparse reward often makes learning inefficient and unstable. Hand-crafted dense reward can partially fix this,
but it is expensive to design, difficult to scale, and brittle when the task or environment changes.
In parallel, retrieval-augmented systems have shown value across many domains whenever the agent needs
external knowledge to handle rare or out-of-distribution situations.

At the same time, robotics research has moved strongly along three relevant directions:
\begin{enumerate}
  \item retrieval/memory for embodied agents and planning \citep{liu2024embodiedrag, li2025embodiedrag, xu2024raea};
  \item retrieval from demonstrations or human videos for control and generalization \citep{johns2024dinobot, huang2025strap, jain2024vid2robot};
  \item reward modeling with vision-language models to reduce dependence on manual reward design \citep{ma2024vlmfeedback, yang2025rgvlm, song2026roboreward}.
\end{enumerate}

The still-open gap lies at the intersection of those directions:
\textbf{using multimodal retrieval to enrich reward signals for continuous manipulation},
especially under OOD and long-horizon conditions. That is where this proposal places its contribution.

\section{Objectives, Research Questions, and Hypotheses}

\subsection{Objectives}

\begin{itemize}
  \item Build a retrieval-conditioned reward modeling pipeline for robot manipulation in simulation.
  \item Measure whether retrieval improves reward quality and policy learning relative to reward models without retrieval.
  \item Evaluate how robust this approach is on long-horizon tasks and under OOD conditions.
\end{itemize}

\subsection{Research Questions}

\begin{enumerate}
  \item Does context retrieved from video demonstrations and instruction/manual text help the reward model estimate task progress more accurately?
  \item Does retrieval-conditioned reward improve \textit{success rate}, \textit{sample efficiency}, and \textit{OOD generalization} compared with no-retrieval baselines?
  \item What retrieval granularity is most effective for reward modeling: full episodes, stage-level chunks, or sub-trajectories?
\end{enumerate}

\subsection{Working Hypotheses}

\begin{itemize}
  \item \textbf{H1}: Multimodal retrieval helps the reward model distinguish meaningful intermediate progress from misleading states.
  \item \textbf{H2}: Stage-level or sub-trajectory retrieval is more effective than episode-level retrieval because it reduces context noise.
  \item \textbf{H3}: The benefit of retrieval is larger on OOD splits and long-horizon tasks than on short, fully in-distribution tasks.
\end{itemize}
